\documentclass[10pt]{article}
\usepackage[margin=1in]{geometry}
\usepackage{booktabs,graphicx,hyperref,amsmath}
\usepackage[numbers]{natbib}
\graphicspath{{../docs/figures/}}
\title{VITRINE: Explainable Static PE Triage and a Correlational Account of Temporal Drift\\from EMBER 2018 to EMBER2024}
\author{Rakshit\\VIT Chennai}
\date{Preprint, October 2026 (not peer reviewed)}
\begin{document}
\maketitle

\begin{abstract}
Static machine-learning malware detectors are usually reported as a single score on a single,
temporally adjacent test set. We present VITRINE, a static-only triage pipeline that scores Windows PE
files with gradient-boosted trees over 91 named features, explains each verdict with exact TreeSHAP,
tags import-level capabilities and synthesizes benign-validated YARA rules, and we use it to measure
temporal drift and rank candidate sources for it. Trained on EMBER 2018 and tested on EMBER2024, the model
loses @@DROPAUC@@ ROC AUC (0.9914 to 0.9730) and, at its own calibrated threshold, recall falls by
@@DROPREC@@ points (88.0\% to 58.6\%) while the false-positive rate \emph{drops} to 0.07\%: the decay is
silent, and the EMBER authors' own 2018 model decays the same way. A drift-weighted ranking
(PSI $\times$ mean$|$SHAP$|$; correlational) over features that are defined identically in both
schemas and that LIEF~0.9 and pefile extract identically on 1{,}069 benign files points to
@@RANKSUMMARY@@. The top raw-PSI feature (embedded MZ count) is a schema-definition change, not drift,
and four section-entropy features differ between the two extractors on identical files. Removing the
ranked features does not recover robustness, and an oracle intervention that gives the ranked features
their 2018 distributions @@ALIGNSUMMARY@@. On identical test rows the published EMBER 2018
model is ahead of VITRINE by @@PUBAUC@@ AUC. All experiments use pre-extracted features; no malware
binaries are handled.
\end{abstract}

\section{Introduction}
Analysts need more than \texttt{predict() = 0.87}: they need reasons, a candidate signature, and an
idea of whether the model still works on today's files. EMBER \citep{anderson2018ember} established
static-feature baselines; EMBER2024 \citep{joyce2025ember2024} renewed the data; TESSERACT
\citep{pendlebury2019tesseract} showed that evaluation must respect time. We combine these: an
interpretable triage model whose features have names, evaluated across two dataset generations six
years apart, with a correlational feature-level drift ranking. Because the two generations were
collected differently and extracted by different parsers (LIEF~0.9 for EMBER 2018, pefile for
EMBER2024), we treat every named shift as a hypothesis and test it three ways.

Contributions: (i) a cross-time evaluation EMBER 2018 $\leftrightarrow$ EMBER2024 with seed-pooled
bootstrap intervals for every reported statistic, a same-size control, a schema-translation control
and the published EMBER 2018 model as a reference; (ii) a cross-generation \emph{definition and
extractor parity audit} of named features: both feature definitions, and both extractors, run on the
same benign files; (iii) a correlational drift-weighted ranking with bootstrap rank intervals,
binning and family-composition controls, a negative removal ablation and an oracle marginal-alignment
intervention; (iv) a re-scoring of both datasets' published reference models on our exact rows, so
the gap to the state of the art is measured rather than assumed; (v) a benign-validated structural
YARA synthesizer with abstention, compared with like-for-like benign-filtered baselines.

\section{Related work}
Static PE detectors with hashed features and gradient-boosted trees
\citep{anderson2018ember,ke2017lightgbm,emberrepo} are strong but opaque; EMBER 2018 is described in a
CAMLIS talk rather than the original paper \citep{roth2019camlis}. SHAP \citep{lundberg2017shap} and exact
TreeSHAP \citep{lundberg2020tree} give additive attributions for tree ensembles \citep{chen2016xgboost}.
Automatic rule generation ranges from frequency-based tools such as yarGen \citep{yargen} to
biclustering \citep{raff2020autoyara}; capa \citep{capa} detects capabilities from code.

\textbf{Drift.} TESSERACT \citep{pendlebury2019tesseract} formalised time-aware splits and the AUT
metric. Transcend \citep{jordaney2017transcend} and its revision \citep{barbero2022transcending} detect
drifting samples with conformal evaluation and reject them. CADE \citep{yang2021cade} detects drifting
samples with contrastive learning and explains them by the features that move a sample furthest from its
class. Drift forensics \citep{chow2023drift} analyses which features and families drive the
performance drop of malware classifiers over time, and \citet{kalny2026drift} quantify per-family
drift on EMBER2024 with rule-based representations and compare against Transcend. Our setting is
narrower: two dataset \emph{generations} whose feature schemas and PE parsers differ. What we add is the
parity audit that shows which named features can be compared across the generations at all, and the
evidence that a plausible-looking ranking is confirmed neither by removal nor by intervention.

\section{Method}
\textbf{Parsing and features.} A bounds-checked pure-Python PE parser emits the EMBER v2 raw schema;
91 named features (header flags, mitigations, section entropy, import groups, string counters,
capability indicators) are computed from it. EMBER2024 feature-version-3 records are translated to v2
(section names lower-cased, a few counters approximated).

\textbf{Model and explanation.} XGBoost (600 trees, depth 10, learning rate 0.08); thresholds at 1\%
and 0.1\% FPR calibrated on a held-out period; exact TreeSHAP contributions with sentence templates.
SHAP is faithful in the deletion sense: replacing a malware sample's top-$k$ features with the benign
median lowers its score @@SHAPRATIO@@ times more than replacing $k$ random features.

\textbf{Rules.} Structural atoms (imports, section names, imphash) kept only if rare in a training
benign pool; the smallest $N$-of threshold with zero benign hits, else abstain.

\textbf{Definition and extractor parity.} Each named feature is classed by how its v2 and v3 values are
defined: \emph{redefined} (4 string counters; v3 has only the nearest regex), \emph{proxy} (6 flags
derived from other fields), \emph{parser} or \emph{byte-level}. For redefined counters both regex
definitions are run on the same 600 System32 files. For parser features, elastic/ember's v2 extractor
with LIEF~0.9.0 and thrember's v3 extractor with pefile are run on the same 1{,}069 benign DLL, PYD and
EXE files unpacked from PyPI Windows wheels inside a GitHub Actions runner; a parser feature whose two
values agree within 1\% on fewer than 99\% of files is classed \emph{extractor-sensitive}.
\emph{Faithful} features are the parser (extractor-consistent) and byte-level ones.

\textbf{Drift ranking.} Per feature and class, the population stability index (PSI) between the 2018
and 2024 test sets on point-mass-aware bins (every value holding $\ge$10\% of the reference gets its
own bin; the rest are reference deciles; shares floored at $10^{-4}$), weighted by the 2018 model's
mean $|$SHAP$|$: $w = \max(\mathrm{PSI}_\mathrm{benign}, \mathrm{PSI}_\mathrm{malware}) \times
\overline{|\phi|}$. Classes and periods are resampled jointly over features (200 bootstrap resamples)
for PSI and rank intervals. Controls: the earlier reference-decile PSI, pooled-quantile bins, and the
top 2018 malware family (xtrat, 19.3\% of the 2018 test malware, absent from 2024) removed. The
ranking is correlational.

\textbf{Interventions.} \emph{Removal}: retrain the 2018 model without the top-10 faithful features
(3 seeds) and compare with 100 group-matched random removals, each seed-matched; Monte Carlo
$p = (1 + \#\text{as extreme}) / (1 + 100)$. \emph{Alignment}: on the frozen 2018 models, map each
2024 test value of a feature group through its within-class empirical CDF onto the 2018 test
quantiles of the same class, keeping within-class ranks, and re-score. Because the mapping uses labels
it is an oracle counterfactual for analysis, not a correction.

\section{Experimental setup}
EMBER 2018: deterministic 50\% sha256-prefix subsample of the 600k labelled training rows (memory
limit of a 16\,GB laptop), 275{,}732 rows to fit, 2018-10 for calibration, 200{,}000 test files.
EMBER2024 Win32: first 8\,MiB of each of 64 weekly files, a sha256-range sample of about 7\% of the
paper's split (135{,}423 files; 43.5\% of the test subsample is malicious, against 50\% in the paper),
weeks 0--47 fit, 48--51 calibrate, 52--63 test (27{,}715 files). Three seeds per trained model;
1{,}000-resample stratified bootstrap pooled over seeds; paired differences on shared resamples; the
2018-test minus 2024-test drop resamples the two test sets independently. Every result file records
the commit, command and, for Actions jobs, the run id.

\section{Results}
\begin{table}[h]\centering\small
\begin{tabular}{lccc}\toprule
Train $\rightarrow$ test & ROC AUC (95\% CI) & TPR@1\% FPR (95\% CI) & FPR / recall at source threshold (95\% CI)\\\midrule
@@TABLEROWS@@
\bottomrule
\end{tabular}
\caption{Cross-time results, 91 named features, XGBoost, 3 seeds; intervals are seed-pooled stratified
bootstraps. The published EMBER 2018 model is one fixed model scored by us.}\label{tab:cross}\end{table}

\textbf{Silent decay.} The same 2018 models lose @@DROPAUCCI@@ AUC and @@DROPRECCI@@ points of recall
at their own threshold from the 2018 to the 2024 test set, while their false-positive rate falls to
0.07\%. At equal training size an in-period 2024 model beats the 2018 model on 2024 by
@@PAIREDAUC@@ AUC. The 2018 model's weekly AUC over the 64 EMBER2024 weeks ranges @@WEEKRANGE@@
(Figure~\ref{fig:drift}). The silent decay is not specific to named features: the EMBER authors'
hashed-feature 2018 model drops @@PUBDROP@@ AUC to 0.9833 (0.9823--0.9843) and, at the threshold that
gives it 1\% FPR on the 2018 test set, keeps 68.5\% recall (67.7--69.4\%) at an FPR of 0.006\% (one
false positive in 15{,}672 benign files). Schema translation costs only 0.0006 AUC (0.0003--0.0008) when
the EMBER2024 paper configuration is trained on native vs translated vectors of identical rows; that
control cannot detect a feature defined differently in 2018 and 2024.

\begin{figure}[t]\centering
\includegraphics[width=\linewidth]{drift_ember2024.png}
\caption{Left: AUC of the 2018-trained model on each EMBER2024 week (95\% bootstrap band). Right:
drift-weighted ranking of faithful features with 95\% bootstrap intervals; the top raw-PSI feature is
shown greyed out because its rank is a definition change.}\label{fig:drift}\end{figure}

\textbf{Parity audit.} \texttt{n\_embedded\_mz} counts the bytes \texttt{MZ} anywhere in a file in v2
and strings containing \texttt{!This program} in v3. On 600 System32 files (Windows Server 2022, in GitHub Actions) the two definitions agree
exactly on @@MZAGREE@@ of files (total variation @@MZTV@@); on the 1{,}069 parity files the extractors
agree on 0.1\%. Its cross-dataset PSI is therefore a definition change. Among parser features, header,
import, export and size features agree on at least 99\% of the parity files, but the four
section-entropy features do not (within-1\% agreement 13--82\%; extractor-only PSI up to
@@ENTPSI@@), so they leave the faithful set. The parity corpus is benign open-source software, not
EMBER's files; agreement there is necessary, not sufficient.

\textbf{Ranking.} @@RANKTEXT@@

\textbf{Removal and alignment.} @@ABLTEXT@@ @@ALIGNTEXT@@

\textbf{Reference models.} The EMBER authors' published 2018 model, scored by us on the same 200k test
rows, reaches AUC 0.9964 (0.9962--0.9966), 96.5\% / 87.0\% TPR at 1\% / 0.1\% FPR (their notebook:
86.8\% at 0.1\% FPR). Seed-pooled over three seeds, VITRINE is behind by @@PUBDIFF@@. The gap mixes data
(the published model saw about twice the labelled rows and 1{,}000 trees) and representation, which we
do not separate. Under the published EMBER 2017 setup (feature v1, LightGBM defaults, run in GitHub
Actions, run 37005292317) we obtain AUC 0.99911 (0.99905--0.99917), 98.2\% and 93.0\% TPR at 1\% /
0.1\% FPR against the paper's 0.99911, 98.2\% and 92.99\%. The released EMBER2024 Win32 model scores
AUC 0.9983 (0.9980--0.9985) on our test subsample (paper, full split: 0.9984); our retrain of its
configuration on 107{,}708 rows (about 7\% of the paper's) does not reach that (0.9966, 0.9961--0.9970).

\textbf{Rules.} Over 15 families out of time, VITRINE's rules reach a mean per-family coverage of
16.8\% (family-bootstrap 95\% CI 2.1--35.0\%) with 5 false positives per 100k benign files, against
13.9\% and 7 for an imphash baseline filtered on the same benign pool. The difference, +2.9 points
($-11.0$ to $+20.2$), is not significant (exact sign-flip $p = 0.75$) and comes from one family
(xtrat); pooled over all 63{,}389 test siblings the rules cover 38.0\%, 78\% of it xtrat.

\textbf{Robustness.} In a feature-space attack combining overlay, section, import and header edits,
detection falls to 36.6--44.1\% for all score models (Wilson intervals about $\pm$1.8 points); an
import-based capability floor appears to help (58.7\%), but with donors that trip no high-risk
capability it gives 47.1\% (45.3--48.9\%), and it costs 11.7\% benign FPR (11.5--11.9\%).

\section{Limitations}
EMBER 2018 training uses half of the labelled rows, so the gap to the published model mixes data and
representation; EMBER2024 is a 7\% subsample that is not class-balanced. The drift ranking is a
correlational PSI $\times$ mean$|$SHAP$|$ ranking. Parser-based features were extracted by LIEF~0.9
(EMBER 2018) and pefile/thrember (EMBER2024); our parity run compares the two extractors on benign
open-source DLLs, not on EMBER's own files or on malformed and packed malware, where parsers diverge
more, so the named toolchain shifts may still partly reflect extractor differences. Ranks depend on
binning and on family composition (Section~5). The alignment intervention uses labels and keeps the
2024 dependence between features, so aligned rows are feature combinations the 2018 models never saw;
its negative result rules out a marginal-only shift, not a role for the ranked features. Adversarial
edits are simulated in feature space. Capability tags were compared with capa only as agreement, not
ground truth.

\section{Ethics and reproducibility}
No malware binaries were downloaded, stored or executed; all data are pre-extracted features
(EMBER 2018: MIT; EMBER2024: Apache-2.0), read-only scans of local benign Windows files, and benign
open-source DLLs from PyPI parsed inside a disposable GitHub Actions runner. Code, result files with
per-file provenance (commit, command, Actions run id), per-range SHA-256 manifests and exact commands
are at \url{https://github.com/rakshit-737/vitrine-malware-triage}.

\section*{Acknowledgements}
Code, experiments and this text were developed with AI assistance (Claude, Anthropic); the author
reviewed and is responsible for the content.

\bibliographystyle{plainnat}
\bibliography{refs}
\end{document}
